\section{结论}

本文针对固定参数 CLAHE 在低照度增强中存在的噪声放大和块效应问题，提出了一种噪声抑制的局部方差自适应 CLAHE 方法。该方法通过残差域 MAD 噪声估计构造结构方差图，并基于分位数归一化生成 tile 级自适应 clip map，再利用 tile 间双线性融合减弱块边界突变。实验结果表明，本文方法更适合被定位为一种可解释的折中型改进：其在亮度提升、结构恢复和噪声控制之间形成了一定平衡，而融合策略对降低块效应具有明确作用。

同时，本文也验证了方法的局限性：当前参数和实现下，方法并未在所有定量指标上优于固定 CLAHE 或去噪后 CLAHE；真实图上仍存在偏灰和局部条带倾向；在更强退化场景中，虽然噪声放大相对可控，但结构恢复仍有提升空间。后续工作可从更稳定的无参考自然度建模、更精细的 tile 插值方式以及颜色恢复机制三个方向继续改进。
